\section{Original Architecture Based on LTE Mode 3}

Before the migration to 5G, the \texttt{plexe\_hetnet} subproject used as its third interface a cellular component derived from the LTE architecture already present in PLEXE subprojects dedicated to C-V2X. Within the multi-technology system, its role was to complement 802.11p and VLC as a redundant cellular channel for the dissemination of cooperative beacons.

From an architectural point of view, this component followed the SimuLTE model: vehicles modeled as \emph{user equipment}, \emph{eNodeB} base stations, a simplified network core, and the global modules required for radio management \cite{virdis2014simulte}. In particular, the \emph{eNodeB} is the LTE entity responsible for radio resource allocation, while the \emph{Binder} provides a global view of the system state that is useful for interference estimation and radio medium coordination. In the scenario used in \texttt{plexe\_hetnet}, two \emph{eNodeB}s covered the road segment traversed by the platoon and were connected to each other so as to support vehicle mobility and \emph{handover} transients.

The central point for the project is that cellular communication did not take place as traditional application traffic routed toward a remote server, but as one-to-many vehicular communication over the C-V2X component. In 3GPP terms, the configuration corresponds to \emph{Mode 3}: vehicles operate under cellular network coverage and sidelink resource management remains assisted by the infrastructure \cite{3gpp-21914-1400}. This distinguishes the LTE component both from 802.11p, where medium access is distributed and contention-based, and from cellular approaches in which resource selection is carried out entirely autonomously on board the vehicle.

In the concrete operation of \texttt{plexe\_hetnet}, each vehicle was therefore equipped with an LTE NIC with D2D support and associated with a cell, with the possibility of dynamic association and \emph{handover} while moving. The beacons generated by the platooning application were forwarded to the LTE radio driver and then transmitted over the C-V2X component with one-to-many dissemination semantics, thus preserving the same redundancy logic already used for 802.11p and VLC. In other words, the application continued to see the cellular network as a third radio available for sending cooperative messages, while the details of LTE communication remained delegated to the underlying stack.

This setup had two important consequences. The first is that the cellular channel introduced characteristics into the multi-technology system that differed from those of the other two interfaces: broader coverage and coordinated access, but also infrastructure dependence and the presence of transients related to cell changes. The second is that the higher-level logic of \emph{SafeSwitch} could treat LTE as an additional interface to be monitored through PDR, without needing to know the handover logic or the implementation details: from the point of view of the platooning protocol, LTE simply remained one of the three possible sources of received beacons.
